\chapter{The Hypercomputational Fibration}
\label{sec:fibration}
% ===================================================================

\section{Background: Turing's Oracle Hierarchy}

Turing~\cite{turing1939} described a transfinite tower of
computational systems above the Turing-computable.
At the base is the class of Turing-computable functions.
At the first stage above is the class of functions computable
relative to the \emph{halting oracle} $\emptyset'$: the oracle that
decides, for any Turing machine $M$ and input $x$, whether $M$
halts on $x$.
At the next stage is the class computable relative to
$\emptyset''$, the halting oracle for machines with access to
$\emptyset'$.
The tower continues through all countable ordinals and beyond.

Each stage is strictly more powerful than the stage below:
there are functions computable at stage $\alpha+1$ that are not
computable at stage $\alpha$, and this gap cannot be closed by
any finite amount of additional computation at stage $\alpha$.

\section{The Subcategory of Turing-Complete GSLTs}

\begin{definition}[Turing-Complete GSLT]
A GSLT $\GSLT$ is \emph{Turing-complete} if the class of functions
computable by terms of $\GSLT$ (via some standard encoding of
natural numbers) coincides with the class of Turing-computable
functions.
The full subcategory of Turing-complete GSLTs is written
$\mathbf{GSLT}_{\mathrm{TC}} \subseteq \mathbf{GSLT}$.
\end{definition}

\begin{example}
The $\lambda$-calculus, the $\pi$-calculus, and the Rho calculus
are all Turing-complete.
The category $\mathbf{GSLT}_{\mathrm{TC}}$ is therefore large and
contains the primary examples of interest.
\end{example}

\section{Oracle Extension of a GSLT}

\begin{definition}[Oracle Term]
Let $\GSLT \in \mathbf{GSLT}_{\mathrm{TC}}$.
An \emph{oracle term} for $\GSLT$ is a distinguished term
$\oracle \in \terms(\GSLT^{(1)})$ in an extension of $\GSLT$,
equipped with rewrite rules of the form:
\[
  \oracle \inter \lceil (M, x) \rceil \;\rewrite\;
  \begin{cases}
    \lceil \mathtt{halt} \rceil & \text{if } M \text{ halts on } x \\
    \lceil \mathtt{loop} \rceil & \text{otherwise}
  \end{cases}
\]
where $\lceil \cdot \rceil$ denotes the encoding of Turing machine
descriptions and inputs as terms of $\GSLT$, and $\inter$ is the
interaction operator of $\GSLT$.
The oracle answers halting queries in one interaction step.
\end{definition}

\begin{remark}[The Oracle as a Communicating Process]
In the Rho calculus, the oracle is naturally expressed as a
persistent receive on a distinguished channel $x_{\oracle}$:
\[
  \oracle \;=\; \mathtt{for}(q \leftarrow x_\oracle)\,
  \bigl( \text{compute halting of } {*}q \text{ and reply} \bigr)
\]
This makes the oracle a \emph{communicating process} rather than
a magic tape: it participates in the same interaction mechanism as
every other term.
This is essential for the fibration to be a fibration of GSLTs
rather than a fibration that escapes the GSLT framework at each
oracle step.
\end{remark}

\begin{definition}[Oracle Extension]
\label{def:oracle-ext}
Let $\GSLT \in \mathbf{GSLT}_{\mathrm{TC}}$.
The \emph{first oracle extension} $\fiber{1}$ is the GSLT obtained
from $\GSLT$ by:
\begin{enumerate}[label=(\roman*)]
  \item extending the grammar with the oracle term $\oracle$;
  \item adding the oracle rewrite rules to $\rules$;
  \item extending the equations $\eqs$ to account for oracle
        interactions.
\end{enumerate}
The \emph{$\alpha$-th oracle extension} $\fiber{\alpha}$ is defined
by transfinite induction:
\begin{itemize}[itemsep=2pt]
  \item $\fiber{0} = \GSLT$;
  \item $\fiber{\alpha+1}$ is the oracle extension of $\fiber{\alpha}$,
        with oracle deciding halting in $\fiber{\alpha}$;
  \item $\fiber{\lambda} = \varinjlim_{\alpha < \lambda} \fiber{\alpha}$
        for limit ordinals $\lambda$ (colimit in $\mathbf{GSLT}$).
\end{itemize}
\end{definition}

\section{The Fibration}

\begin{definition}[Hypercomputational Fibration]
The \emph{hypercomputational fibration} over $\mathbf{GSLT}_{\mathrm{TC}}$
is the functor
\[
  \Phi : \mathbf{GSLT}_{\mathrm{TC}} \times \Ord \;\to\; \mathbf{GSLT}
\]
sending $(\GSLT, \alpha)$ to $\fiber{\alpha}$, together with:
\begin{itemize}[itemsep=2pt]
  \item projection morphisms $\proj{\alpha} : \fiber{\alpha+1} \to \fiber{\alpha}$
        in $\mathbf{GSLT}$ for each ordinal $\alpha$;
  \item inclusion morphisms $\iota_\alpha : \fiber{\alpha} \to \fiber{\alpha+1}$
        embedding each stage into the next.
\end{itemize}
The fiber over $\GSLT$ is the tower
$\fiber{0} \hookrightarrow \fiber{1} \hookrightarrow \fiber{2}
\hookrightarrow \cdots$
\end{definition}

\begin{proposition}[Projections are Morphisms]
Each projection $\proj{\alpha} : \fiber{\alpha+1} \to \fiber{\alpha}$
is a morphism in $\mathbf{GSLT}$: it preserves bisimulation.
If $P \bisim_{\alpha+1} Q$ in $\fiber{\alpha+1}$, then
$\proj{\alpha}(P) \bisim_\alpha \proj{\alpha}(Q)$ in $\fiber{\alpha}$.
\end{proposition}

\begin{proof}[Proof sketch]
Any bisimulation in $\fiber{\alpha+1}$ that uses only oracle-free
transitions is already a bisimulation in $\fiber{\alpha}$.
The projection discards oracle interactions; the remaining
bisimulation relation is preserved.
The details require checking that the oracle rewrite rules are
projected consistently with the extension of $\eqs$.
\end{proof}

\begin{remark}[Projections are Lossy]
The projection $\proj{\alpha}$ is in general \emph{not} a
bisimulation equivalence: there exist $P, Q$ with
$P \bisim_{\alpha+1} Q$ but $P \not\bisim_\alpha Q$.
Going down the fibration collapses distinctions.
Going up reveals them.
\end{remark}

% ===================================================================
